\documentclass[12pt]{article}
\usepackage[margin=1in]{geometry}
\usepackage{amsmath}
\usepackage{graphicx}
\usepackage{booktabs}
\usepackage{hyperref}
\usepackage{cite}
\usepackage{setspace}
\usepackage{caption}
\usepackage{float}
\usepackage{listings}
\usepackage{xcolor}
\usepackage{microtype}

\onehalfspacing

\title{Beyond Energy Consumption: Carbon Emission as an Evaluation Metric\\
for Contextual Bandit-Based LLM Inference Routing}

\author{Yuna Jung\\
Murty Sunak Quantitative Computing Lab\\
Claremont McKenna College\\
Summer Research Program 2026\\
\texttt{yjung29@cmc.edu}}

\date{July 2026}

\begin{document}

\maketitle

% ─────────────────────────────────────────────────────────────────────────────
\begin{abstract}
This study investigates whether carbon emission should be introduced as an
evaluation metric for environmentally-friendly large language model (LLM)
inference routing algorithms. We replicate GreenServ \cite{ziller2026}, a
contextual multi-armed bandit routing framework that dynamically assigns
queries to models from a pool of 16 language models to optimize the
accuracy-energy tradeoff, on the Laguna high-performance computing cluster
at USC/CARC. GreenServ currently evaluates routing performance using mean
normalized accuracy, GPU energy consumption (Wh), and latency, but does not
measure carbon emissions. We extend GreenServ by integrating CodeCarbon
\cite{codecarbon2023}, an open-source Python library that measures actual
energy consumption of GPU, CPU, and RAM and converts it to equivalent CO$_2$
(gCO$_2$) using regional grid carbon intensity. We conduct two experiments:
(1) a per-model carbon benchmark measuring carbon emissions for 15 of
GreenServ's 16 language models across 10 test queries, finding that carbon
scales with parameter count but is also influenced by model architecture; and
(2) a live routing experiment running all 8 routing strategies across
GreenServ's full 2,500-query benchmark dataset with CodeCarbon measuring
carbon per routing decision. Our goal is to produce an Accuracy
vs.\ Carbon Emissions Pareto plot comparable to GreenServ's existing Accuracy
vs.\ Energy plot, and to determine whether energy consumption is a sufficient
proxy for carbon emissions or whether carbon must be measured independently
to accurately evaluate the environmental impact of LLM routing algorithms.
\end{abstract}

% ─────────────────────────────────────────────────────────────────────────────
\section{Introduction}

The rapid growth of large language model (LLM) deployment has raised
significant concerns about the environmental cost of AI inference. LLMs are
power-intensive: a single inference call requires matrix multiplications
across billions of parameters, drawing hundreds of watts from GPU hardware.
At the scale of millions of daily queries across commercial deployments, the
cumulative energy and carbon footprint becomes substantial. A 2024 TED talk
by AI researcher Sasha Luccioni highlighted the environmental impact of
machine learning data centers, motivating the question at the center of this
work: can we design smarter routing systems that direct queries to
appropriately-sized models, reducing environmental impact without sacrificing
accuracy?

This question led to an exploration of Small Language Models (SLMs),
generally defined as models in the 1--9 billion parameter range
\cite{gupta2026}, as more efficient alternatives to large models for tasks
that do not require full frontier-model capability. The fundamental
observation is that many real-world queries are answerable by smaller, more
efficient models --- routing them to a large model is wasteful.

During a literature review of existing routing approaches, we found that
while several frameworks address computational or monetary efficiency
\cite{chen2023frugalgpt, routellm2024}, few explicitly optimize for
environmental metrics. GreenServ \cite{ziller2026} is a notable exception:
it uses a contextual multi-armed bandit (MAB) algorithm --- specifically
LinUCB \cite{li2010contextual} --- to route queries across a pool of 16
models, optimizing a scalarized reward combining normalized accuracy and GPU
energy consumption. However, GreenServ measures environmental cost only
through GPU electricity usage in watt-hours (Wh), and does not account for
carbon emissions.

This distinction matters because the carbon footprint of a given amount of
electricity depends heavily on the energy mix of the local grid. The same
watt-hour of compute in a coal-powered region emits far more CO$_2$ than in
a renewables-heavy grid such as California's. Moreover, carbon emissions,
water usage, and electricity are three distinct environmental impacts of data
center operations \cite{luccioni2023}, and optimizing for one does not
necessarily optimize for the others.

This report presents two contributions: (1) a per-model carbon measurement
study using CodeCarbon \cite{codecarbon2023} on GreenServ's 16 models, and
(2) an integration of CodeCarbon into GreenServ's live inference routing
pipeline to measure carbon emissions per routing decision across all 8
routing strategies on the full 2,500-query benchmark. The central research
question is: \textit{Is GPU energy consumption a sufficient proxy for carbon
emissions when evaluating the environmental friendliness of LLM inference
routing algorithms?}

% ─────────────────────────────────────────────────────────────────────────────
\section{Background and Related Work}

\subsection{Environmental Costs of Language Models}

The environmental impact of LLMs spans three primary dimensions: electricity
consumption, carbon emissions, and water usage \cite{luccioni2023,
strubell2019}. Training large models requires enormous compute --- for
example, training GPT-3 consumed an estimated 1,287 MWh of electricity
\cite{patterson2021}. While inference per query is much cheaper than
training, its cumulative cost at scale is significant. Luccioni et al.\
\cite{luccioni2023hf} benchmarked the carbon footprint of multiple models
and tasks using CodeCarbon, finding substantial variance across architectures
and hardware. Schwartz et al.\ \cite{schwartz2020green} introduced the
concept of ``Green AI,'' arguing that efficiency should be treated as a
first-class research objective alongside accuracy.

\subsection{LLM and SLM Routing}

Routing frameworks aim to reduce inference cost by directing queries to
models of appropriate capacity. FrugalGPT \cite{chen2023frugalgpt} uses a
cascade approach: queries are sent to progressively larger models until a
confidence threshold is met, optimizing for API cost. RouteLLM
\cite{routellm2024} trains a classifier on human preference data to predict
whether a query needs a strong or weak model. Both frameworks use cost in
dollars or token budgets as proxies for resource consumption, not direct
energy or carbon measurements.

\subsection{GreenServ}

GreenServ \cite{ziller2026} is the closest prior work to this study. It
implements a context-aware routing framework using LinUCB, a contextual
bandit algorithm that assumes a linear relationship between query context and
model reward. The system extracts three query features: task type (MMLU,
GSM8K, HellaSwag, WinoGrande, CNN/DailyMail), semantic cluster (via online
K-means on sentence embeddings), and text complexity (Flesch Reading Ease
score). These features are one-hot encoded into a context vector of dimension
$d = 12$, which is used to select among 16 models spanning 0.5B to 34B
parameters. The reward function is:
\begin{equation}
r_t = (1 - \lambda) \cdot \text{Acc}_\text{norm} - \lambda \cdot E_\text{norm}
\end{equation}
where $\lambda \in [0,1]$ balances accuracy and energy efficiency.
GreenServ evaluates routing algorithms on mean normalized accuracy and total
energy consumption (Wh), measured via the Zeus GPU monitoring library
\cite{zeus2023}. It does not measure carbon emissions.

\subsection{CodeCarbon}

CodeCarbon \cite{codecarbon2023} is an open-source Python library that
measures the carbon footprint of computational workloads. It monitors GPU,
CPU, and RAM power consumption via hardware interfaces (NVIDIA NVML for GPU,
Intel RAPL for CPU) and converts total energy (kWh) to CO$_2$ equivalent
(kgCO$_2$eq) using the carbon intensity of the local electricity grid,
sourced from regional databases or the ElectricityMap API. Two tracker
modes are available: \texttt{EmissionsTracker} (online, uses real-time grid
data) and \texttt{OfflineEmissionsTracker} (offline, uses a static carbon
intensity database by country). Prior work \cite{luccioni2023hf} used
CodeCarbon to benchmark the carbon footprint of HuggingFace models across
tasks and hardware, demonstrating its suitability for measuring inference
carbon emissions.

% ─────────────────────────────────────────────────────────────────────────────
\section{Methods}

\subsection{Infrastructure}

All experiments were conducted on the Laguna high-performance computing
cluster at the University of Southern California's Center for Advanced
Research Computing (CARC), accessed by Claremont McKenna College. The
cluster uses SLURM for job scheduling. Compute nodes used in this project
were equipped with dual NVIDIA L40S GPUs (46 GB VRAM each, 92 GB total) and
AMD EPYC 9354 32-core processors with 128 GB RAM. Model weights were stored
in project storage (\texttt{/project/JehoPark\_1895/Yuna\_SLM/}), which
provides 338 TB of shared network-attached storage.

Since Laguna does not support Docker for security reasons, GreenServ's
PostgreSQL database dependency was migrated to Apptainer (formerly
Singularity), the HPC-standard container runtime. The migration required
pulling the official PostgreSQL 15 Docker image via Apptainer, initializing
the database directory in project storage, and creating a writable Unix
socket directory to bypass the container's read-only filesystem. This
migration is documented in full in the lab notebook for reproducibility.

\subsection{GreenServ Replication}

We cloned GreenServ's public repository
(\url{https://github.com/TZData1/llm-inference-router}) and set up a Python
3.11 virtual environment with the provided \texttt{requirements.txt}. The
PostgreSQL database was populated by loading \texttt{db/llm\_db\_backup.sql}
into a local instance, yielding 40,000 pregenerated inference results (2,500
queries $\times$ 16 models $\times$ multiple runs). All four GreenServ
experiments were executed via the provided Makefile targets
(\texttt{make exp01-full} through \texttt{make exp04-full}), using SLURM
for job management with \texttt{nohup} to survive session disconnections.

The key infrastructure challenge was that GreenServ's experiments replay
pregenerated results from the database rather than running live inference.
This means neither Zeus nor CodeCarbon fires during the standard experiment
runs. Carbon measurements for the routing algorithm comparison therefore
required a separate live inference pipeline (Section~\ref{sec:live}).

\subsection{Experiment 1: Per-Model Carbon Benchmark}
\label{sec:permodel}

To establish a baseline understanding of how carbon emissions vary across
GreenServ's model pool, we ran each model through 10 test queries and
measured carbon emissions using CodeCarbon's
\texttt{OfflineEmissionsTracker} with \texttt{country\_iso\_code="USA"}.

The 10 queries were general knowledge prompts (e.g., ``What is the capital
of France?'', ``What is machine learning?'') designed to elicit short,
factual responses. For each model, we:
\begin{enumerate}
  \item Loaded the model weights from the HuggingFace cache in project
  storage using \texttt{AutoModelForCausalLM} in \texttt{bfloat16}
  precision with \texttt{device\_map="auto"}.
  \item Started a CodeCarbon \texttt{OfflineEmissionsTracker} tracker.
  \item Ran inference on all 10 queries with \texttt{max\_new\_tokens=50}.
  \item Stopped the tracker and recorded total emissions in gCO$_2$.
  \item Deleted the model and called \texttt{torch.cuda.empty\_cache()} to
  free GPU memory before the next model.
\end{enumerate}

Of GreenServ's 16 models, 15 were successfully measured. Yi-34B
(\texttt{01-ai/Yi-34B-Chat}) initially failed due to transformer
compatibility issues but was confirmed accessible on HuggingFace and is
included in subsequent experiments. Results were saved to
\texttt{live\_carbon\_results.csv}.

\subsection{Experiment 2: Live Routing Carbon Measurement}
\label{sec:live}

To measure carbon emissions \textit{per routing algorithm} --- rather than
per model --- we designed a live inference experiment that mirrors
GreenServ's evaluation setup but replaces the pregenerated database replay
with actual model inference, with CodeCarbon measuring carbon during each
inference call.

\subsubsection{Routing Algorithms}

We evaluated all 8 routing strategies tested in GreenServ:
\begin{itemize}
  \item \textbf{LinUCB (C)}: GreenServ's primary algorithm, contextual
  bandit with upper confidence bound exploration ($\alpha = 0.1$,
  $\lambda_\text{reg} = 0.05$).
  \item \textbf{Thompson Sampling (C)}: Bayesian posterior sampling
  ($\sigma = 0.25$, prior variance = 2).
  \item \textbf{$\varepsilon$-Greedy (C)}: Contextual epsilon-greedy
  ($\varepsilon_0 = 1.0$, $\delta = 0.98$, $\varepsilon_\text{min} = 0.01$).
  \item \textbf{$\varepsilon$-Greedy (NC)}: Non-contextual epsilon-greedy
  (same hyperparameters, $\lambda = 0.25$).
  \item \textbf{Random}: Uniformly random model selection.
  \item \textbf{Largest}: Always selects Yi-34B (largest model).
  \item \textbf{Smallest}: Always selects Qwen2.5-0.5B (smallest model).
  \item \textbf{Accuracy}: Always selects Gemma-3-27B (highest accuracy
  model).
\end{itemize}

\subsubsection{Feature Extraction}

We used GreenServ's full feature extraction pipeline via
\texttt{FeatureService}, which computes:
\begin{itemize}
  \item \textbf{Task type}: Derived directly from the source dataset label
  (MMLU, GSM8K, HellaSwag, WinoGrande, CNN/DailyMail), one-hot encoded.
  \item \textbf{Semantic cluster}: Online K-means ($k=3$) on
  \texttt{all-MiniLM-L6-v2} sentence embeddings, one-hot encoded.
  \item \textbf{Text complexity}: Flesch Reading Ease score discretized
  into 3 bins, one-hot encoded.
  \item \textbf{Intercept}: A constant 1.0 appended to the context vector.
\end{itemize}
This yields a context vector of dimension $d = 12$, matching GreenServ's
original setup.

\subsubsection{Reward Computation}

To evaluate routing decisions without re-running model evaluation (which
would require ground-truth labels and evaluation metrics), we used
GreenServ's pregenerated accuracy and energy values from the database as the
reward signal, consistent with GreenServ's offline evaluation approach. The
reward for selecting model $m$ for query $q$ is:
\begin{equation}
r(m, q) = (1 - \lambda) \cdot \text{Acc}_\text{norm}(m, q) - \lambda \cdot E_\text{norm}(m, q)
\end{equation}
with $\lambda = 0.5$ (equal weight on accuracy and energy efficiency).

\subsubsection{Carbon Measurement}

For each algorithm, we wrapped the full 2,500-query inference run in a
single CodeCarbon \texttt{OfflineEmissionsTracker} instance. This captures
the total GPU, CPU, and RAM energy consumed by that algorithm's routing
decisions --- which implicitly reflects model selection frequency. When an
algorithm like LinUCB primarily routes to small models, it consumes less GPU
power; when the Largest baseline always calls Yi-34B, GPU power draw is
consistently high. CodeCarbon captures these differences at the hardware
level.

Note on multi-GPU usage: experiments were submitted with
\texttt{--gres=gpu:2} to provide 92 GB total VRAM across two L40S GPUs.
The \texttt{device\_map="balanced"} setting in \texttt{from\_pretrained()}
distributes model layers evenly across both GPUs. This reduces model loading
time for large models and prevents out-of-memory errors between algorithm
runs.

\subsubsection{GPU Memory Management}

A key implementation challenge was GPU out-of-memory errors between
algorithms. Models loaded during one algorithm's run remained in GPU memory
and caused the next algorithm's model loads to fail. We resolved this by
explicitly calling \texttt{del model}, \texttt{del tokenizer},
\texttt{torch.cuda.empty\_cache()}, and \texttt{gc.collect()} after each
algorithm's run before starting the next.

\subsubsection{Errors Encountered}

Table~\ref{tab:errors} summarizes the engineering challenges encountered
during development of the live routing experiment.

\begin{table}[H]
\centering
\caption{Errors encountered during live routing experiment development}
\label{tab:errors}
\begin{tabular}{p{3.5cm} p{5.5cm} p{5cm}}
\toprule
\textbf{Error} & \textbf{Cause} & \textbf{Fix} \\
\midrule
\texttt{TypeError}: missing \texttt{model\_ids}, \texttt{context\_dimension} &
Bandit constructors require model list and context dim as first positional args &
Added \texttt{model\_ids} and \texttt{context\_dimension} to all bandit instantiations \\
\addlinespace
\texttt{AttributeError}: no attribute \texttt{select} &
GreenServ's method is \texttt{select\_model()}, not \texttt{select()} &
Changed to \texttt{algo.select\_model(context)} \\
\addlinespace
Gated repo error: \texttt{aya-expanse-32b} &
\texttt{models.yaml} contains extra models beyond GreenServ's 16 &
Hardcoded \texttt{GREENSERV\_MODEL\_IDS} list of 16 models \\
\addlinespace
\texttt{CUDA OutOfMemoryError} between algorithms &
Loaded models remained in GPU memory across algorithm runs &
Added explicit GPU cleanup (\texttt{del}, \texttt{empty\_cache()}, \texttt{gc.collect()}) \\
\addlinespace
\texttt{RuntimeError}: layer not mapped to device &
\texttt{device\_map="auto"} failed to assign all layers across 2 GPUs &
Changed to \texttt{device\_map="balanced"} and \texttt{use\_cache=False} \\
\addlinespace
\texttt{sbatch} node config unavailable &
Requested 8 GPUs but nodes have max 2 L40S GPUs &
Changed \texttt{--gres=gpu:8} to \texttt{--gres=gpu:2} \\
\bottomrule
\end{tabular}
\end{table}

% ─────────────────────────────────────────────────────────────────────────────
\section{Results}

\subsection{Experiment 1: Per-Model Carbon Benchmark}

Table~\ref{tab:permodel} shows carbon emissions per model for 10 queries
each, measured on the Laguna L40S GPU using
\texttt{OfflineEmissionsTracker}.

\begin{table}[H]
\centering
\caption{Carbon emissions per model (10 queries, Laguna L40S GPU)}
\label{tab:permodel}
\begin{tabular}{lrr}
\toprule
\textbf{Model} & \textbf{Parameters} & \textbf{Carbon (gCO$_2$)} \\
\midrule
Qwen2.5-0.5B  & 0.5B  & 0.171 \\
Qwen2.5-1.5B  & 1.5B  & 0.207 \\
Phi-4-mini-4B & 3.8B  & 0.213 \\
Qwen2.5-3B    & 3B    & 0.221 \\
Qwen2.5-7B    & 7B    & 0.355 \\
Mistral-7B    & 7B    & 0.373 \\
Gemma-3-1B    & 1B    & 0.634 \\
Qwen2.5-14B   & 14B   & 0.646 \\
Phi-4-14B     & 14B   & 0.925 \\
Gemma-3-4B    & 4B    & 1.234 \\
Llama-3.1-8B  & 8B    & 1.807 \\
Llama-3.2-1B  & 1B    & [pending] \\
Llama-3.2-3B  & 3B    & [pending] \\
Gemma-3-12B   & 12B   & 2.529 \\
Gemma-3-27B   & 27B   & 4.470 \\
Yi-34B        & 34B   & [pending] \\
\bottomrule
\end{tabular}
\end{table}

Two key findings emerge. First, carbon emissions generally increase with
parameter count. Second, model architecture significantly affects carbon
efficiency: Phi-4-mini (3.8B parameters) emits less carbon per query than
Qwen2.5-1.5B (1.5B parameters), suggesting that architectural efficiency
differences outweigh raw parameter counts in some cases. Similarly, Gemma-3
models emit substantially more carbon than Qwen2.5 models of similar size.

\subsection{Experiment 2: Live Routing Carbon Measurement}

The live routing experiment was submitted to Laguna's GPU partition multiple
times across several sessions but did not complete within the timeframe of
this research project. The experiment encountered a fundamental infrastructure
limitation: \textbf{model loading overhead}.

In GreenServ's original design, all model inference results are pregenerated
and stored in a PostgreSQL database. During experiments, the router replays
these results without ever loading model weights. This makes GreenServ's
experiments fast and GPU-memory-efficient. Our live routing experiment
departed from this design by actually loading and running each model during
the routing loop --- which is necessary to measure real CodeCarbon emissions
but introduces significant overhead.

Specifically, when a routing algorithm (e.g.\ LinUCB) selects a different
model for each query, the system must load that model's weights from disk
into GPU memory before inference can occur. Each model load takes 30--120
seconds depending on model size. With 15 models and 2,500 queries, a routing
algorithm that frequently switches between models could trigger hundreds of
model loads, making the experiment infeasible within a single SLURM job
window.

Table~\ref{tab:exp2errors} summarizes the specific errors encountered across
multiple submission attempts.

\begin{table}[H]
\centering
\caption{Errors encountered during Experiment 2 execution}
\label{tab:exp2errors}
\begin{tabular}{p{4cm} p{5cm} p{4.5cm}}
\toprule
\textbf{Error} & \textbf{Cause} & \textbf{Fix Applied} \\
\midrule
Job stuck at Query 0/2500 for 2+ hours &
Yi-34B (34B parameters) barely fit in 92GB VRAM across 2 L40S GPUs;
inference crawled due to memory pressure and offloading &
Removed Yi-34B from model pool \\
\addlinespace
\texttt{RuntimeError}: layer not mapped to device &
\texttt{device\_map="auto"} failed to distribute layers across 2 GPUs &
Changed to \texttt{device\_map="balanced"}, added \texttt{use\_cache=False} \\
\addlinespace
Job still stuck at Query 0/2500 after removing Yi-34B &
LinUCB routes to many different models in early queries;
each model switch requires a full reload from disk &
Fundamental design limitation --- not resolved within project timeframe \\
\bottomrule
\end{tabular}
\end{table}

\subsubsection{Why This Experiment Cannot Run As Designed}

The root cause is a mismatch between GreenServ's pregenerated-replay
architecture and the requirements of live CodeCarbon measurement. To produce
an Accuracy vs.\ Carbon Emissions Pareto plot equivalent to GreenServ's
Accuracy vs.\ Energy plot, one of two approaches is needed:

\begin{enumerate}
\item \textbf{Preload all models simultaneously}: Load all 15 models into
GPU memory at once and keep them resident throughout the experiment, avoiding
reload overhead. This requires approximately 150--200 GB of VRAM for the
full model pool, exceeding the 92 GB available on Laguna's dual-L40S nodes.

\item \textbf{Two-phase carbon attribution}: Run each model separately on
all 2,500 queries with CodeCarbon (collecting per-model, per-query carbon
measurements), then use GreenServ's pregenerated routing decisions to
attribute carbon to each algorithm post-hoc. This avoids live model
switching but requires completing the per-model inference job (job 71861,
still running at time of writing).
\end{enumerate}

The two-phase approach is the most feasible path forward given available
hardware. Job 71861 (per-model live inference) had completed 4 of 15 models
at the time this report was finalized. Future work should complete this job
and implement the carbon attribution post-processing step.

% ─────────────────────────────────────────────────────────────────────────────
\section{Discussion}

\subsection{Is Energy a Sufficient Proxy for Carbon?}

The central research question of this study is whether GPU energy
consumption (Wh), as measured by Zeus in GreenServ, is a sufficient proxy
for carbon emissions (gCO$_2$), as measured by CodeCarbon. There are two
possible outcomes:

\textbf{If energy and carbon are linearly correlated across algorithms:}
Energy consumption is a sufficient proxy. The relative ranking of algorithms
by environmental impact is the same whether measured in Wh or gCO$_2$.
GreenServ's existing metric captures environmental cost accurately, and
adding CodeCarbon provides only a unit conversion. This would still be
useful for communicating impact in carbon terms (e.g., ``LinUCB saves X
grams of CO$_2$ per 1,000 queries''), but would not change the conclusions
about which algorithms are most environmentally friendly.

\textbf{If energy and carbon diverge across algorithms:}
Carbon is a necessary independent metric. This could occur because CodeCarbon
measures CPU and RAM energy in addition to GPU energy (which Zeus does not),
and because different routing decisions may produce different CPU-to-GPU
energy ratios. For example, feature extraction (CPU-bound) runs before every
inference call, and its cost varies with query complexity regardless of which
model is ultimately selected. If some algorithms trigger more complex feature
extraction paths, their CPU energy overhead may not be captured by Zeus's
GPU-only measurement.

\subsection{Per-Model Results: Architecture Matters}

The per-model carbon benchmark (Experiment 1) reveals that parameter count
alone does not determine carbon footprint. The Phi-4-mini model (3.8B
parameters) emits less carbon than Qwen2.5-1.5B (1.5B parameters) per 10
queries. This suggests that architectural improvements --- such as more
efficient attention mechanisms, better training data, or improved
tokenization --- can yield models that are both smaller and more carbon
efficient. This finding has implications for model pool selection in routing
frameworks: a router that selects models based on parameter count as a
proxy for energy efficiency may not minimize carbon emissions.

\subsection{Limitations}

\textbf{Inference-only measurement:} Both Zeus and CodeCarbon measure only
inference carbon. The training cost of the 16 pretrained models --- which
may dominate lifetime environmental cost for models used at limited scale
--- is not captured. Faiz et al.\ \cite{faiz2023llmcarbon} provide
framework for estimating training carbon from known compute hours, which
could supplement inference measurements in future work.

\textbf{Offline carbon intensity:} The \texttt{OfflineEmissionsTracker}
uses a static US national average carbon intensity rather than California's
real-time grid data. California's grid is substantially cleaner than the
national average, meaning our carbon estimates are conservative
overestimates. Future work should use the online tracker with
ElectricityMap API for more precise regional measurements.

\textbf{Model caching:} When an algorithm routes to a model already loaded
in GPU memory, no additional loading cost is incurred. However, when it
routes to an unloaded model, a significant loading delay occurs. This means
carbon measurements partially reflect model loading costs in addition to
pure inference costs. GreenServ's original evaluation avoids this by using
pregenerated results; our live experiment does not.

\textbf{Water usage:} Water consumption is the third major environmental
impact of data center operations and is not measured by either Zeus or
CodeCarbon. Future work should incorporate water usage effectiveness (WUE)
metrics.

\textbf{Query type:} GreenServ evaluates only objective benchmark tasks
(multiple choice, math, summarization) with automated accuracy metrics. The
framework's behavior on conversational, subjective, or emotionally sensitive
queries --- where accuracy is harder to define --- remains unstudied.

\subsection{Future Work}

Several directions follow from this study:

\begin{itemize}
  \item \textbf{Online carbon measurement:} Integrate ElectricityMap API
  with \texttt{EmissionsTracker} for real-time California grid carbon
  intensity.
  \item \textbf{Water consumption:} Add water usage as a third environmental
  metric alongside energy and carbon.
  \item \textbf{Lifecycle analysis:} Estimate training carbon for each model
  in the pool using LLMCarbon \cite{faiz2023llmcarbon} and incorporate into
  routing reward.
  \item \textbf{Subjective query routing:} Extend evaluation to
  conversational and open-ended queries where accuracy is assessed via
  LLM-as-judge or human preference.
  \item \textbf{Dynamic model pool:} Investigate criteria for adding or
  removing models from the pool as new releases emerge.
\end{itemize}

% ─────────────────────────────────────────────────────────────────────────────
\section{Conclusion}

This study extends GreenServ, a contextual bandit-based LLM inference
routing framework, by integrating CodeCarbon to measure carbon emissions
alongside the existing GPU energy metric. We demonstrate that carbon
emissions per model scale with parameter count but are also meaningfully
influenced by model architecture --- a finding that has direct implications
for environmentally-conscious model pool design. The live routing experiment
was designed to produce an Accuracy vs.\ Carbon Emissions Pareto plot
comparable to GreenServ's Accuracy vs.\ Energy plot, but could not be
completed within the project timeframe due to GPU memory and model loading
constraints inherent in live inference routing. This limitation itself is a
research finding: integrating real-time carbon measurement into an LLM
routing system requires architectural design choices --- such as preloading
all models or separating per-model carbon collection from routing decision
attribution --- that are not required for energy-only measurement using Zeus.
Future work should complete the two-phase carbon attribution approach and
investigate whether energy consumption and carbon emissions are linearly
correlated across routing algorithms, or whether carbon must be treated as
an independent evaluation metric.

% ─────────────────────────────────────────────────────────────────────────────
\section*{Acknowledgments}

This research was conducted as part of the Claremont McKenna College Summer
Research Program 2026 at the Murty Sunak Quantitative Computing Lab, under
the supervision of Professor Jeho Park. Computing resources were provided by
the USC Center for Advanced Research Computing (CARC) Laguna cluster. This
research uses Claude.ai (Anthropic) as a coding and debugging assistant,
with advisor approval, for writing and debugging Python scripts, navigating
HPC infrastructure, and troubleshooting errors. All research decisions,
methodology, interpretations, and findings are the author's own.

% ─────────────────────────────────────────────────────────────────────────────
\begin{thebibliography}{99}

\bibitem{ziller2026}
Ziller, T., et al.
\textit{GreenServ: Energy-Efficient LLM Inference Routing via Contextual Bandits}.
2026.

\bibitem{codecarbon2023}
Courty, B., et al.
\textit{CodeCarbon: Estimate and Track Carbon Emissions from Machine Learning Computing}.
\url{https://github.com/mlco2/codecarbon}, 2023.

\bibitem{gupta2026}
Gupta, A., et al.
\textit{A Survey of Small Language Models in the 1--9 Billion Parameter Range}.
Amazon Machine Learning, 2026.

\bibitem{li2010contextual}
Li, L., et al.
\textit{A Contextual-Bandit Approach to Personalized News Article Recommendation}.
In \textit{WWW}, 2010.

\bibitem{chen2023frugalgpt}
Chen, L., et al.
\textit{FrugalGPT: How to Use Large Language Models While Reducing Cost and Improving Performance}.
arXiv:2305.05176, 2023.

\bibitem{routellm2024}
Ong, I., et al.
\textit{RouteLLM: Learning to Route LLMs with Preference Data}.
arXiv:2406.18665, 2024.

\bibitem{luccioni2023}
Luccioni, A.S.
\textit{The Environmental Impact of AI}.
TED Talk, 2023.

\bibitem{luccioni2023hf}
Luccioni, A.S., Viguier, S., Ligozat, A.L.
\textit{Estimating the Carbon Footprint of BLOOM, a 176B Parameter Language Model}.
\textit{JMLR}, 2023.

\bibitem{strubell2019}
Strubell, E., Ganesh, A., McCallum, A.
\textit{Energy and Policy Considerations for Deep Learning in NLP}.
In \textit{ACL}, 2019.

\bibitem{patterson2021}
Patterson, D., et al.
\textit{Carbon Emissions and Large Neural Network Training}.
arXiv:2104.10350, 2021.

\bibitem{schwartz2020green}
Schwartz, R., et al.
\textit{Green AI}.
\textit{Communications of the ACM} 63, 12 (2020), 54--63.

\bibitem{zeus2023}
You, J., et al.
\textit{Zeus: Understanding and Optimizing GPU Energy Consumption of DNN Training}.
In \textit{NSDI}, 2023.

\bibitem{faiz2023llmcarbon}
Faiz, A., et al.
\textit{LLMCarbon: Modeling the End-to-End Carbon Footprint of Large Language Models}.
arXiv:2309.14393, 2023.

\end{thebibliography}

\end{document}
